\chapter{Introduction}
\label{chap:intro}

\section{Preliminaries}
\supporting

The paper works with the arboreal representation~$\rho_a$ on the tree of preimages of~$0$; the
formalization works with the Galois groups of the iterates, which are its images on the finite
levels. These are the objects in the statement of Proposition~\ref{prop:criterion}.

\begin{definition}[The iterates~$f_n$]
  \label{def:iterates}
  \lean{QuadraticIterates.iteratedPoly, QuadraticIterates.iteratedPoly_zero,
    QuadraticIterates.iteratedPoly_one, QuadraticIterates.iteratedPoly_succ,
    QuadraticIterates.iteratedPoly_succ_comp, QuadraticIterates.eval_zero_iteratedPoly_succ,
    QuadraticIterates.map_iteratedPoly, QuadraticIterates.aeval_ratCast_iteratedPoly,
    QuadraticIterates.monic_iteratedPoly, QuadraticIterates.natDegree_iteratedPoly,
    QuadraticIterates.iteratedPoly_zero_left,
    QuadraticIterates.ncard_rootSet_iteratedPoly_le}
  \leanok
  $f_0 := X$ and $f_{n+1} := f_n^2 + a$. These are monic of degree~$2^n$, so~$f_n$ has at
  most~$2^n$ roots, and $f_{n+1} = f_n \circ (X^2 + a)$; so $f_n = f^{\circ n}$.
  \formalnote{\lname{iteratedPoly} is defined over any commutative semiring; over~$\QQ$ it is
  written \texttt{f}$\QQ$\texttt{[a, n]}.}
\end{definition}

\begin{definition}[The fields~$K_n$ and the groups~$\Om_n$]
  \label{def:galois}
  \uses{def:iterates}
  \lean{QuadraticIterates.splittingField, QuadraticIterates.GaloisGroup,
    QuadraticIterates.splittingField.instIsSplittingField,
    QuadraticIterates.splittingField.instFiniteDimensional,
    QuadraticIterates.splittingField.instIsGalois,
    QuadraticIterates.splittingField_zero_eq_bot,
    QuadraticIterates.card_galoisGroup_eq_finrank}
  \leanok
  For $a \in \QQ$, $K_n \subseteq \Qbar$ is the subfield generated by the roots of~$f_n$, and
  $\Om_n = \Gal(f_n/\QQ)$. One has $K_0 = \QQ$ and $\#\Om_n = [K_n : \QQ]$.
  \formalnote{$K_n$ is an \lname{IntermediateField} of \lname{AlgebraicClosure ℚ}, so that the
  tower $K_n \subseteq K_{n+1}$ is an inclusion of subfields of one field and relative degrees
  are \lname{IntermediateField.relfinrank}; $\Om_n$ is Mathlib's \lname{Polynomial.Gal}.}
\end{definition}

\begin{definition}[The iterated wreath power~$\Ct{n}$]
  \label{def:wreath}
  \lean{QuadraticIterates.WreathPower, QuadraticIterates.card_wreathPower}
  \leanok
  $\Ct{n}$ is the~$n$-fold iterated regular wreath product of~$C_2$ with itself; it has order
  $2^{2^n - 1}$.
  \formalnote{Mathlib's \lname{IteratedWreathProduct} of \lname{Multiplicative (ZMod 2)}.
  ``$\Om_n \cong \Ct{n}$'' is \lname{Nonempty (GaloisGroup a n ≃* WreathPower n)}, and
  ``$\rho_a$ is surjective'' is the statement that this holds for all~$n$.}
\end{definition}

\section{The question and the sequence \texorpdfstring{$(c_n)$}{(c\_n)}}
\paperitems

Let $a \in \QQ$ be nonzero and $f = f_a = x^2 + a$. If no iterate has multiple roots, the pairs
$(n, \xi)$ with~$\xi$ a root of~$f^{\circ n}$ in a fixed algebraic closure of~$\QQ$ form an infinite
complete binary tree~$T$, in which $(n + 1, \eta)$ is a child of $(n, \xi)$ when $\xi = f(\eta)$.
The absolute Galois group~$\Gal_{\QQ}$ of~$\QQ$ acts on~$T$ by tree automorphisms, which gives the
\emph{arboreal Galois representation} $\rho_a \colon \Gal_{\QQ} \to \operatorname{Aut}(T)$. An
iterate~$f^{\circ n}$ has multiple roots if and only if~$0$ is a root of~$f^{\circ m}$ for some
$1 \le m \le n$; so some iterate has multiple roots if and only if~$0$ is periodic under~$f$,
which for $a \in \QQ$ happens exactly when $a \in \{-1, 0\}$ (Remark~\ref{rem:nonzero}).

\begin{question}
  For which $a \in \QQ$ is~$\rho_a$ surjective?
\end{question}

This is equivalent to the surjectivity, for every~$n$, of the action~$\rho_a^n$ on the finite
tree~$T_n$ of the levels up to~$n$. The automorphism group of~$T_n$ is a Sylow~$2$-subgroup of
the symmetric group on its~$2^n$ leaves, isomorphic to the~$n$th wreath power~$\Ct{n}$ of~$C_2$,
and the image of~$\rho^n_a$ is~$\Om_n$. So~$\rho^n_a$ is surjective if and only if
$\Om_n \cong \Ct{n}$ (Definitions~\ref{def:galois} and~\ref{def:wreath}). The question was first
asked by J.~McKay for $a = 1$; see~\cite{odoni}. Odoni's earlier paper~\cite{odoni1985} sets up
some general theory, and the survey~\cite{jones} describes what is known or expected on this and
related questions.

\begin{definition}[The sequence~$c_n$]
  \label{def:cseq}
  \uses{def:iterates, def:gammaseq}
  \lean{QuadraticIterates.cSeq, QuadraticIterates.cSeq_zero, QuadraticIterates.cSeq_one,
    QuadraticIterates.cSeq_succ, QuadraticIterates.cSeq_two,
    QuadraticIterates.cSeq_succ_eq_neg_one_pow_mul_eval,
    QuadraticIterates.cSeq_succ_eq_neg_one_pow_mul_eval_zero, QuadraticIterates.map_cSeq,
    QuadraticIterates.intCast_cSeq, QuadraticIterates.ratCast_cSeq}
  \leanok
  Define the sequence $(c_n)_{n \ge 0}$ of rational numbers by
  \[ c_0 = 0, \qquad c_1 = -a, \qquad c_{n+1} = f(c_n) = c_n^2 + a \quad \text{for $n \ge 1$.} \]
  This is the orbit of~$0$ under~$f$, \emph{except} that the sign of~$c_1$ is changed (note
  $f(-a) = f(a)$). Thus $c_{n+1} = f_n(a)$ for all~$n$, $c_n = f_n(0)$ for $n \ge 2$, and
  $c_1 = -f_1(0)$.
  \formalnote{\lname{cSeq} is defined for a parameter in any commutative ring, as the
  sequence \lname{gammaSeq} of the even polynomial~$X^2 + a$ with first term~$-a$
  (Definition~\ref{def:gammaseq}); it commutes with ring homomorphisms
  (\lname{map\_cSeq}). This is used with the rings~$\QQ$ and~$\ZZ/N\ZZ$.}
\end{definition}

\begin{remark}[Vanishing terms]
  \label{rem:nonzero}
  $c_n = 0$ for some $n \ge 1$ holds if and only if~$0$ is periodic under~$f$. For $a \in \QQ$,
  this is the case if and only if $a \in \{-1, 0\}$. If~$p$ divides the denominator of~$a$, then
  the~$p$-adic valuation doubles in each step, so~$0$ cannot be periodic; so~$a$ must be an
  integer. For $a \ge 1$ or $a \le -3$ the orbit tends to infinity, leaving $a \in \{-2, -1, 0\}$;
  and for $a = -2$, $0$ is preperiodic ($0 \mapsto -2 \mapsto 2 \mapsto 2$), not periodic.
  \formalnote{The equivalence with $a \in \{0, -1\}$ is \lname{exists\_cSeq\_eq\_zero\_iff}, with
  \lname{cSeq\_ne\_zero} for $a \notin \{-1, 0\}$; both are listed under Lemma~\ref{lem:vp},
  which uses them. The periodicity of~$0$ is not formalized. Most statements do not need the
  remark, as a character with $\chi(c_n) = -1$ already forces $c_n \ne 0$.}
\end{remark}

The~$p$-adic valuations along $(c_n)_{n \ge 1}$ behave as follows.

\begin{lemma}[Rigid divisibility]
  \label{lem:rigid}
  \uses{def:cseq}
  \lean{QuadraticIterates.cSeq_sq_dvd_cSeq_mul_sub, QuadraticIterates.sq_dvd_cSeq_sub,
    QuadraticIterates.cSeq_mul_of_sq_eq_zero, QuadraticIterates.cSeq_eq_self_of_sq_eq_zero,
    QuadraticIterates.cSeq_add_eq_of_sq_eq_zero, QuadraticIterates.cSeq_add_mul_eq_of_sq_eq_zero,
    QuadraticIterates.cSeq_add_eq_of_add_eq, QuadraticIterates.cSeq_add_mul_eq_of_add_eq}
  \leanok
  Let~$p$ be a prime and assume that~$a$ is a~$p$-adic integer. Then~$c_n$ is a~$p$-adic integer
  for all~$n$, and in~$\ZZ_p$ we have $c_{mk} \equiv c_m \bmod c_m^2$ for all $m \ge 2$ and
  $k \ge 1$, and similarly $c_k \equiv -c_1 = a \bmod c_1^2$ for all $k \ge 2$.
  \formalnote{For a parameter in any commutative ring~$R$ (such as~$\ZZ_p$), where the integrality
  is automatic: \lname{cSeq\_sq\_dvd\_cSeq\_mul\_sub} and \lname{sq\_dvd\_cSeq\_sub}. They are the
  images of the ring statements in $R/(c_m^2)$: if $c_m^2 = 0$ ($m \ge 1$), the sequence
  is~$m$-periodic from $n = 2$ on (\lname{cSeq\_add\_eq\_of\_sq\_eq\_zero}), so $c_{mk} = c_m$ for
  $m \ge 2$ (\lname{cSeq\_mul\_of\_sq\_eq\_zero}) and $c_k = a$ for $m = 1$
  (\lname{cSeq\_eq\_self\_of\_sq\_eq\_zero}).}
\end{lemma}
\begin{proof}
  \leanok
  Since $f = x^2 + a$ is even, so is~$f^{\circ n}$ for $n \ge 1$, so
  $f^{\circ n}(z) \equiv f^{\circ n}(0) \bmod z^2$. Fix $m \ge 2$; then $c_m = f^{\circ m}(0)$, and
  by induction on~$k$,
  \[ c_{m(k+1)} = f^{\circ (mk)}(c_m) \equiv f^{\circ (mk)}(0) = c_{mk} \bmod c_m^2. \]
  The case $m = 1$ is similar, using $c_2 = f(c_1) \equiv f(0) = -c_1 \bmod c_1^2$. The
  formalization argues in $\bar R = R/(c_m^2)$, where $\bar c_m^2 = 0$, so
  $\bar c_{m+1} = \bar a$ and $\bar c_{m+2} = \bar a^2 + \bar a = \bar c_2$; periodicity then
  propagates along the recursion from $n = 2$ on.
\end{proof}

\begin{lemma}[The~$p$-adic valuations along~$(c_n)$]
  \label{lem:vp}
  \uses{lem:rigid, def:wseq, lem:reduction, lem:wseq-res, lem:wseq-mul, lem:signs}
  \lean{QuadraticIterates.padicValRat_cSeq_trichotomy,
    QuadraticIterates.padicValRat_cSeq_of_padicValRat_neg,
    QuadraticIterates.padicValRat_cSeq_of_dvd_left, QuadraticIterates.padicValRat_cSeq_eq_ite,
    QuadraticIterates.padicValRat_cSeq_eq_padicValInt, QuadraticIterates.padicValInt_wSeq_eq_ite,
    QuadraticIterates.dvd_wSeq_iff, QuadraticIterates.cSeq_eq_zero_iff_dvd,
    QuadraticIterates.cSeq_ne_zero, QuadraticIterates.exists_cSeq_eq_zero_iff}
  \leanok
  Let~$p$ be a prime and assume that $a \notin \{-1, 0\}$. Then exactly one of the following
  possibilities holds.
  \begin{enumerate}
    \item For all $n \ge 1$, $v_p(c_n) = 0$.
    \item There are $m_p, e_p \in \ZZ_{\ge 1}$ such that $v_p(c_n) = 0$ if $m_p \nmid n$, and
          $v_p(c_n) = e_p$ if $n \ge 1$ and $m_p \mid n$.
    \item There is $e_p \in \ZZ_{< 0}$ such that $v_p(c_n) = 2^{n-1} e_p$ for all $n \ge 1$.
  \end{enumerate}
  \formalnote{\lname{padicValRat\_cSeq\_trichotomy} states that one of the three holds; that at
  most one does is clear from the values at $n = 1$ and $n = m_p$, and is not stated. For
  $a = r/s$, the cases are those of \lname{padicValRat\_cSeq\_of\_padicValRat\_neg} ($p \mid s$,
  $e_p = v_p(a)$), \lname{padicValRat\_cSeq\_of\_dvd\_left} ($p \mid r$, $m_p = 1$,
  $e_p = v_p(r)$) and \lname{padicValRat\_cSeq\_eq\_ite} ($p \nmid rs$, with~$m_p$ the first
  index with $p \mid w_{m_p}$ and $e_p = v_p(w_{m_p})$).}
\end{lemma}
\begin{proof}
  \leanok
  If $v_p(a) < 0$, then $v_p(c_n^2) < v_p(a)$ and so $v_p(c_{n+1}) = 2 v_p(c_n)$; this is the
  third case with $e_p = v_p(a)$. If $v_p(a) \ge 0$, then $v_p(c_n) \ge 0$ for all~$n$
  (Lemma~\ref{lem:rigid}). If $v_p(c_n) = 0$ for all $n \ge 1$, we are in the first case.
  Otherwise let $m_p = \min\{m \ge 1 : v_p(c_m) > 0\}$ and $e_p = v_p(c_{m_p})$: the sequence
  modulo~$p$, starting from~$0$, is periodic with period~$m_p$, and Lemma~\ref{lem:rigid} with
  $m = m_p$ gives $v_p(c_{k m_p}) = e_p$. Compare the proof of~\cite{paper}, Lemma~1.1(b).

  The formalization writes $a = r/s$ and $c_n = r w_n / s^{2^{n-1}}$ (Lemma~\ref{lem:reduction}).
  For $p \nmid s$, $v_p(c_n) = v_p(r w_n)$, and~$w_n$ is prime to~$r$ and to~$s$
  (Lemma~\ref{lem:wseq-res}); this settles $p \mid r$. For $p \nmid rs$, modulo~$p$ the vanishing
  of~$\bar c_n$ is that of~$w_n$; if~$m$ is the first index with $\bar c_m = 0$, then
  $\bar c_n = 0$ if and only if $m \mid n$ (\lname{cSeq\_eq\_zero\_iff\_dvd}), and
  $w_{mk} = w_m T$ with $T \equiv s^{2^{mk-1} - 2^{m-1}} \bmod w_m$ prime to~$p$
  (Lemma~\ref{lem:wseq-mul}). The terms are nonzero by Remark~\ref{rem:nonzero}: for
  non-integral~$a$ by the third case, for integral $a \notin \{-1, 0\}$ because $a > 0$ or
  $a \le -2$ and Lemma~\ref{lem:signs} applies.
\end{proof}

\begin{definition}[$2$-independence]
  \label{def:twoindep}
  \lean{QuadraticIterates.TwoIndependent, QuadraticIterates.TwoIndependent.ne_zero,
    QuadraticIterates.TwoIndependent.not_isSquare,
    QuadraticIterates.twoIndependent_iff_linearIndependent,
    QuadraticIterates.TwoIndependent.of_snoc, QuadraticIterates.twoIndependent_snoc_iff,
    QuadraticIterates.TwoIndependent.comp, QuadraticIterates.twoIndependent_iff_forall_fin,
    QuadraticIterates.twoIndependent_iff_forall_finset,
    Fin.exists_eq_map_castSuccEmb_or_insert_last, sqClass}
  \leanok
  A family $(r_i)_{i \in I}$ of rational numbers is \emph{$2$-independent} if for every finite
  subset $S \subseteq I$, $\prod_{i \in S} r_i$ is a square only when $S = \emptyset$. Otherwise
  the family is \emph{$2$-dependent}. (When all~$r_i$ are nonzero, this says that the images of
  the~$r_i$ in $\QQ^\times/\QQ^{\times 2}$ are linearly independent over~$\FF_2$.) A family
  containing zero is~$2$-dependent.
  \formalnote{\lname{TwoIndependent} is stated for a family in any commutative monoid; the
  translation into linear independence of the classes \lname{sqClass} in $L^\times/L^{\times 2}$
  holds for any field~$L$. The family $c_1, \dots, c_n$ is $i \mapsto c_{i+1}$ on \lname{Fin n}.}
\end{definition}

The following is the equivalence ``a) $\Leftrightarrow$ b)'' in part~1) of the Theorem in
Section~1 of~\cite{paper}.

\begin{proposition}[Stoll's criterion]
  \label{prop:criterion}
  \uses{def:galois, def:wreath, def:twoindep, def:cseq}
  \lean{QuadraticIterates.nonempty_mulEquiv_iff_twoIndependent,
    QuadraticIterates.irreducible_iteratedPoly_of_twoIndependent,
    QuadraticIterates.forall_nonempty_mulEquiv_iff_twoIndependent}
  \leanok
  Fix $a \in \QQ$ and $n \in \ZZ_{\ge 1}$. Then~$\rho^n_a$ is defined and surjective if and only
  if $(c_1, \dots, c_n)$ is~$2$-independent. Consequently, $\rho_a$ is defined and surjective if
  and only if the family $(c_n)_{n \ge 1}$ is~$2$-independent.
  \formalnote{Stated as $\Om_n \cong \Ct{n}$ if and only if $c_1, \dots, c_n$
  are~$2$-independent, for every $n \ge 0$ (\lname{nonempty\_mulEquiv\_iff\_twoIndependent}), and
  for all~$n$ at once (\lname{forall\_nonempty\_mulEquiv\_iff\_twoIndependent}, with the family
  \lname{fun n : ℕ ↦ cSeq a (n + 1)}).}
\end{proposition}
\begin{proof}
  \leanok
  \uses{lem:criterion-irr, lem:irred, lem:irred-of-iso}
  The proof is the subject of Appendix~\ref{chap:criterion}. By
  Lemma~\ref{lem:criterion-irr} it suffices that $f_0, \dots, f_n$ are irreducible under either
  hypothesis. If $c_1, \dots, c_n$ are~$2$-independent, none of them is a square, and
  Lemma~\ref{lem:irred} applies; if $\Om_n \cong \Ct{n}$, Lemma~\ref{lem:irred-of-iso} does. For
  the second statement, a finite subset of $\{1, 2, \dots\}$ lies in some $\{1, \dots, n\}$.
\end{proof}

\cite{paper} defines~$2$-independence only for families of nonzero numbers and assumes that~$-a$
is not a square (which implies that all~$c_n$ are nonzero). The statement remains true in the
general form above: when~$-a$ is a square, $\rho^1_a$ is not surjective, hence neither
is~$\rho^n_a$ for $n \ge 1$, and $c_1 = -a$ is a square, so $(c_1, \dots, c_n)$ is~$2$-dependent.
This includes $a = 0$ and $a = -1$, where~$\rho_a$ is not defined. So the approach is to show
that $(c_n)_{n \ge 1}$ is~$2$-independent for $a \in \QQ$ with suitable properties.

The proofs in~\cite{paper, li2020, li2021} follow~\cite{odoni} in considering instead the
sequence~$(b_n)$ obtained from~$(c_n)$ by multiplicative Möbius inversion. When
$a \in \ZZ \setminus \{-1, 0\}$, $(c_n)$ is a rigid divisibility sequence in~$\ZZ$
(Lemma~\ref{lem:vp}), so the~$b_n$ are pairwise coprime integers; $(c_n)_{n \ge 1}$
is~$2$-independent if and only if $(b_n)_{n \ge 1}$ is, and this follows if~$|b_n|$ is never a
square for $n \ge 2$ and~$b_1$ is not a square. That gives a simpler criterion --- single terms
instead of arbitrary finite products --- at the price of more complicated objects. The paper
shows that the results of~\cite{paper, li2020, li2021} are obtained as easily by working
directly with~$(c_n)$, and extends them to rational parameters. Its main result
(Theorem~\ref{thm:main}) is that for essentially half of all rational~$a$ outside the interval
$(-2, \alpha)$, with~$\alpha$ the real root of $x^3 + 2x^2 + 2x + 2$, $\rho_a$ is surjective as
soon as (and therefore if and only if) $\rho_a^2$ is, that is, as soon as none of~$-a$,
$a^2 + a$ and $-a(a^2 + a)$ is a square.
